\documentclass[11pt]{article}
% ============================================================================
% preamble.tex -- shared style for every document in docs/
%
% Each document starts with:
%     \documentclass[11pt]{article}
%     % ============================================================================
% preamble.tex -- shared style for every document in docs/
%
% Each document starts with:
%     \documentclass[11pt]{article}
%     % ============================================================================
% preamble.tex -- shared style for every document in docs/
%
% Each document starts with:
%     \documentclass[11pt]{article}
%     \input{preamble}
%     \newcommand{\doctitle}{...}
%
% Team names live in ONE place (below). Edit them here and every document
% updates on the next build.
% ============================================================================

% ---------------------------------------------------------------- team
\newcommand{\projname}{FPGA int8 MLP Accelerator}
\newcommand{\repourl}{https://github.com/itsKannapi/fpga-mlp-accelerator}
\newcommand{\memberA}{Joey (\texttt{itsKannapi})}
\newcommand{\memberB}{Teammate~2 (TBD)}
\newcommand{\memberC}{Teammate~3 (TBD)}

% ---------------------------------------------------------------- packages
\usepackage[margin=1in]{geometry}
\usepackage[T1]{fontenc}
\usepackage{lmodern}
\usepackage{microtype}
\usepackage{amsmath}
\usepackage{booktabs}
\usepackage{tabularx}
\usepackage{longtable}
\usepackage{array}
\usepackage{enumitem}
\usepackage{xcolor}
\usepackage{listings}
\usepackage{tikz}
\usetikzlibrary{arrows.meta, positioning, fit, calc}
\usepackage[most]{tcolorbox}
\usepackage{fancyhdr}
\usepackage{titlesec}
\usepackage{needspace}
\usepackage[hidelinks]{hyperref}
\hypersetup{colorlinks=true, linkcolor=accent, urlcolor=accent, citecolor=accent}

% ---------------------------------------------------------------- colors
\definecolor{accent}{HTML}{1F4E79}
\definecolor{muted}{HTML}{5F6B7A}
\definecolor{rule}{HTML}{C9D1DA}
\definecolor{okc}{HTML}{2E7D32}
\definecolor{progc}{HTML}{B26A00}
\definecolor{todoc}{HTML}{8A8F98}
\definecolor{codebg}{HTML}{F5F7FA}

% ---------------------------------------------------------------- layout
\setcounter{tocdepth}{1}   % sections only, keeps the contents to a few lines
\setlength{\parindent}{0pt}
\setlength{\parskip}{0.55em}
\setlist{itemsep=0.15em, topsep=0.3em}
\renewcommand{\arraystretch}{1.18}

\titleformat{\section}{\Large\bfseries\color{accent}}{\thesection}{0.7em}{}
\titleformat{\subsection}{\large\bfseries}{\thesubsection}{0.6em}{}
\titleformat{\subsubsection}{\normalsize\bfseries}{\thesubsubsection}{0.5em}{}
\titlespacing*{\section}{0pt}{1.4em}{0.5em}

\pagestyle{fancy}
\fancyhf{}
\fancyhead[L]{\small\color{muted}\projname}
\fancyhead[R]{\small\color{muted}\doctitle}
\fancyfoot[C]{\small\color{muted}\thepage}
\renewcommand{\headrulewidth}{0.4pt}
\renewcommand{\headrule}{\color{rule}\hrule width\headwidth height\headrulewidth}

% ---------------------------------------------------------------- title block
\newcommand{\makedoctitle}[2]{%  #1 = subtitle, #2 = last-updated date
  \thispagestyle{plain}
  {\color{muted}\small\MakeUppercase{\projname}}\par\vspace{0.2em}
  {\Huge\bfseries\color{accent}\doctitle}\par\vspace{0.4em}
  {\large #1}\par\vspace{0.6em}
  {\small\color{muted}Team: \memberA, \memberB, \memberC\quad$\cdot$\quad
   Last updated #2\quad$\cdot$\quad\href{\repourl}{\texttt{\repourl}}}\par
  \vspace{0.4em}{\color{rule}\rule{\linewidth}{0.6pt}}\par\vspace{0.6em}
}

% ---------------------------------------------------------------- status tags
\newcommand{\tagbox}[2]{\tcbox[on line, colback=#1!12, colframe=#1!12,
  coltext=#1!70!black, boxsep=0pt, left=3pt, right=3pt, top=1.5pt, bottom=1.5pt,
  arc=2pt]{\footnotesize\bfseries #2}}
\newcommand{\done}{\tagbox{okc}{DONE}}
\newcommand{\inprog}{\tagbox{progc}{IN PROGRESS}}
\newcommand{\todo}{\tagbox{todoc}{TO DO}}
\newcommand{\stretchgoal}{\tagbox{accent}{STRETCH}}

% ---------------------------------------------------------------- callouts
\newtcolorbox{note}[1][Note]{colback=accent!5, colframe=accent!60,
  boxrule=0.6pt, arc=2pt, left=6pt, right=6pt, top=4pt, bottom=4pt,
  fonttitle=\bfseries\small, title=#1, coltitle=white}
\newtcolorbox{warn}[1][Watch out]{colback=progc!6, colframe=progc!70,
  boxrule=0.6pt, arc=2pt, left=6pt, right=6pt, top=4pt, bottom=4pt,
  fonttitle=\bfseries\small, title=#1, coltitle=white}

% ---------------------------------------------------------------- code
\lstdefinelanguage{SV}{
  morekeywords={module,endmodule,input,output,wire,reg,signed,always,posedge,
    negedge,if,else,begin,end,assign,parameter,localparam,generate,endgenerate,
    for,genvar,case,endcase,function,endfunction,integer,initial},
  sensitive=true, morecomment=[l]{//}, morecomment=[s]{/*}{*/},
  morestring=[b]", morestring=[b]`}
\lstset{
  basicstyle=\ttfamily\small, backgroundcolor=\color{codebg},
  keywordstyle=\color{accent}\bfseries, commentstyle=\color{muted}\itshape,
  frame=single, rulecolor=\color{rule}, framesep=4pt, xleftmargin=4pt,
  xrightmargin=4pt, columns=fullflexible, keepspaces=true, breaklines=true,
  showstringspaces=false, upquote=true, aboveskip=0.6em, belowskip=0.6em}
\newcommand{\cmd}[1]{\texttt{#1}}
\newcommand{\file}[1]{\texttt{#1}}

%     \newcommand{\doctitle}{...}
%
% Team names live in ONE place (below). Edit them here and every document
% updates on the next build.
% ============================================================================

% ---------------------------------------------------------------- team
\newcommand{\projname}{FPGA int8 MLP Accelerator}
\newcommand{\repourl}{https://github.com/itsKannapi/fpga-mlp-accelerator}
\newcommand{\memberA}{Joey (\texttt{itsKannapi})}
\newcommand{\memberB}{Teammate~2 (TBD)}
\newcommand{\memberC}{Teammate~3 (TBD)}

% ---------------------------------------------------------------- packages
\usepackage[margin=1in]{geometry}
\usepackage[T1]{fontenc}
\usepackage{lmodern}
\usepackage{microtype}
\usepackage{amsmath}
\usepackage{booktabs}
\usepackage{tabularx}
\usepackage{longtable}
\usepackage{array}
\usepackage{enumitem}
\usepackage{xcolor}
\usepackage{listings}
\usepackage{tikz}
\usetikzlibrary{arrows.meta, positioning, fit, calc}
\usepackage[most]{tcolorbox}
\usepackage{fancyhdr}
\usepackage{titlesec}
\usepackage{needspace}
\usepackage[hidelinks]{hyperref}
\hypersetup{colorlinks=true, linkcolor=accent, urlcolor=accent, citecolor=accent}

% ---------------------------------------------------------------- colors
\definecolor{accent}{HTML}{1F4E79}
\definecolor{muted}{HTML}{5F6B7A}
\definecolor{rule}{HTML}{C9D1DA}
\definecolor{okc}{HTML}{2E7D32}
\definecolor{progc}{HTML}{B26A00}
\definecolor{todoc}{HTML}{8A8F98}
\definecolor{codebg}{HTML}{F5F7FA}

% ---------------------------------------------------------------- layout
\setcounter{tocdepth}{1}   % sections only, keeps the contents to a few lines
\setlength{\parindent}{0pt}
\setlength{\parskip}{0.55em}
\setlist{itemsep=0.15em, topsep=0.3em}
\renewcommand{\arraystretch}{1.18}

\titleformat{\section}{\Large\bfseries\color{accent}}{\thesection}{0.7em}{}
\titleformat{\subsection}{\large\bfseries}{\thesubsection}{0.6em}{}
\titleformat{\subsubsection}{\normalsize\bfseries}{\thesubsubsection}{0.5em}{}
\titlespacing*{\section}{0pt}{1.4em}{0.5em}

\pagestyle{fancy}
\fancyhf{}
\fancyhead[L]{\small\color{muted}\projname}
\fancyhead[R]{\small\color{muted}\doctitle}
\fancyfoot[C]{\small\color{muted}\thepage}
\renewcommand{\headrulewidth}{0.4pt}
\renewcommand{\headrule}{\color{rule}\hrule width\headwidth height\headrulewidth}

% ---------------------------------------------------------------- title block
\newcommand{\makedoctitle}[2]{%  #1 = subtitle, #2 = last-updated date
  \thispagestyle{plain}
  {\color{muted}\small\MakeUppercase{\projname}}\par\vspace{0.2em}
  {\Huge\bfseries\color{accent}\doctitle}\par\vspace{0.4em}
  {\large #1}\par\vspace{0.6em}
  {\small\color{muted}Team: \memberA, \memberB, \memberC\quad$\cdot$\quad
   Last updated #2\quad$\cdot$\quad\href{\repourl}{\texttt{\repourl}}}\par
  \vspace{0.4em}{\color{rule}\rule{\linewidth}{0.6pt}}\par\vspace{0.6em}
}

% ---------------------------------------------------------------- status tags
\newcommand{\tagbox}[2]{\tcbox[on line, colback=#1!12, colframe=#1!12,
  coltext=#1!70!black, boxsep=0pt, left=3pt, right=3pt, top=1.5pt, bottom=1.5pt,
  arc=2pt]{\footnotesize\bfseries #2}}
\newcommand{\done}{\tagbox{okc}{DONE}}
\newcommand{\inprog}{\tagbox{progc}{IN PROGRESS}}
\newcommand{\todo}{\tagbox{todoc}{TO DO}}
\newcommand{\stretchgoal}{\tagbox{accent}{STRETCH}}

% ---------------------------------------------------------------- callouts
\newtcolorbox{note}[1][Note]{colback=accent!5, colframe=accent!60,
  boxrule=0.6pt, arc=2pt, left=6pt, right=6pt, top=4pt, bottom=4pt,
  fonttitle=\bfseries\small, title=#1, coltitle=white}
\newtcolorbox{warn}[1][Watch out]{colback=progc!6, colframe=progc!70,
  boxrule=0.6pt, arc=2pt, left=6pt, right=6pt, top=4pt, bottom=4pt,
  fonttitle=\bfseries\small, title=#1, coltitle=white}

% ---------------------------------------------------------------- code
\lstdefinelanguage{SV}{
  morekeywords={module,endmodule,input,output,wire,reg,signed,always,posedge,
    negedge,if,else,begin,end,assign,parameter,localparam,generate,endgenerate,
    for,genvar,case,endcase,function,endfunction,integer,initial},
  sensitive=true, morecomment=[l]{//}, morecomment=[s]{/*}{*/},
  morestring=[b]", morestring=[b]`}
\lstset{
  basicstyle=\ttfamily\small, backgroundcolor=\color{codebg},
  keywordstyle=\color{accent}\bfseries, commentstyle=\color{muted}\itshape,
  frame=single, rulecolor=\color{rule}, framesep=4pt, xleftmargin=4pt,
  xrightmargin=4pt, columns=fullflexible, keepspaces=true, breaklines=true,
  showstringspaces=false, upquote=true, aboveskip=0.6em, belowskip=0.6em}
\newcommand{\cmd}[1]{\texttt{#1}}
\newcommand{\file}[1]{\texttt{#1}}

%     \newcommand{\doctitle}{...}
%
% Team names live in ONE place (below). Edit them here and every document
% updates on the next build.
% ============================================================================

% ---------------------------------------------------------------- team
\newcommand{\projname}{FPGA int8 MLP Accelerator}
\newcommand{\repourl}{https://github.com/itsKannapi/fpga-mlp-accelerator}
\newcommand{\memberA}{Joey (\texttt{itsKannapi})}
\newcommand{\memberB}{Teammate~2 (TBD)}
\newcommand{\memberC}{Teammate~3 (TBD)}

% ---------------------------------------------------------------- packages
\usepackage[margin=1in]{geometry}
\usepackage[T1]{fontenc}
\usepackage{lmodern}
\usepackage{microtype}
\usepackage{amsmath}
\usepackage{booktabs}
\usepackage{tabularx}
\usepackage{longtable}
\usepackage{array}
\usepackage{enumitem}
\usepackage{xcolor}
\usepackage{listings}
\usepackage{tikz}
\usetikzlibrary{arrows.meta, positioning, fit, calc}
\usepackage[most]{tcolorbox}
\usepackage{fancyhdr}
\usepackage{titlesec}
\usepackage{needspace}
\usepackage[hidelinks]{hyperref}
\hypersetup{colorlinks=true, linkcolor=accent, urlcolor=accent, citecolor=accent}

% ---------------------------------------------------------------- colors
\definecolor{accent}{HTML}{1F4E79}
\definecolor{muted}{HTML}{5F6B7A}
\definecolor{rule}{HTML}{C9D1DA}
\definecolor{okc}{HTML}{2E7D32}
\definecolor{progc}{HTML}{B26A00}
\definecolor{todoc}{HTML}{8A8F98}
\definecolor{codebg}{HTML}{F5F7FA}

% ---------------------------------------------------------------- layout
\setcounter{tocdepth}{1}   % sections only, keeps the contents to a few lines
\setlength{\parindent}{0pt}
\setlength{\parskip}{0.55em}
\setlist{itemsep=0.15em, topsep=0.3em}
\renewcommand{\arraystretch}{1.18}

\titleformat{\section}{\Large\bfseries\color{accent}}{\thesection}{0.7em}{}
\titleformat{\subsection}{\large\bfseries}{\thesubsection}{0.6em}{}
\titleformat{\subsubsection}{\normalsize\bfseries}{\thesubsubsection}{0.5em}{}
\titlespacing*{\section}{0pt}{1.4em}{0.5em}

\pagestyle{fancy}
\fancyhf{}
\fancyhead[L]{\small\color{muted}\projname}
\fancyhead[R]{\small\color{muted}\doctitle}
\fancyfoot[C]{\small\color{muted}\thepage}
\renewcommand{\headrulewidth}{0.4pt}
\renewcommand{\headrule}{\color{rule}\hrule width\headwidth height\headrulewidth}

% ---------------------------------------------------------------- title block
\newcommand{\makedoctitle}[2]{%  #1 = subtitle, #2 = last-updated date
  \thispagestyle{plain}
  {\color{muted}\small\MakeUppercase{\projname}}\par\vspace{0.2em}
  {\Huge\bfseries\color{accent}\doctitle}\par\vspace{0.4em}
  {\large #1}\par\vspace{0.6em}
  {\small\color{muted}Team: \memberA, \memberB, \memberC\quad$\cdot$\quad
   Last updated #2\quad$\cdot$\quad\href{\repourl}{\texttt{\repourl}}}\par
  \vspace{0.4em}{\color{rule}\rule{\linewidth}{0.6pt}}\par\vspace{0.6em}
}

% ---------------------------------------------------------------- status tags
\newcommand{\tagbox}[2]{\tcbox[on line, colback=#1!12, colframe=#1!12,
  coltext=#1!70!black, boxsep=0pt, left=3pt, right=3pt, top=1.5pt, bottom=1.5pt,
  arc=2pt]{\footnotesize\bfseries #2}}
\newcommand{\done}{\tagbox{okc}{DONE}}
\newcommand{\inprog}{\tagbox{progc}{IN PROGRESS}}
\newcommand{\todo}{\tagbox{todoc}{TO DO}}
\newcommand{\stretchgoal}{\tagbox{accent}{STRETCH}}

% ---------------------------------------------------------------- callouts
\newtcolorbox{note}[1][Note]{colback=accent!5, colframe=accent!60,
  boxrule=0.6pt, arc=2pt, left=6pt, right=6pt, top=4pt, bottom=4pt,
  fonttitle=\bfseries\small, title=#1, coltitle=white}
\newtcolorbox{warn}[1][Watch out]{colback=progc!6, colframe=progc!70,
  boxrule=0.6pt, arc=2pt, left=6pt, right=6pt, top=4pt, bottom=4pt,
  fonttitle=\bfseries\small, title=#1, coltitle=white}

% ---------------------------------------------------------------- code
\lstdefinelanguage{SV}{
  morekeywords={module,endmodule,input,output,wire,reg,signed,always,posedge,
    negedge,if,else,begin,end,assign,parameter,localparam,generate,endgenerate,
    for,genvar,case,endcase,function,endfunction,integer,initial},
  sensitive=true, morecomment=[l]{//}, morecomment=[s]{/*}{*/},
  morestring=[b]", morestring=[b]`}
\lstset{
  basicstyle=\ttfamily\small, backgroundcolor=\color{codebg},
  keywordstyle=\color{accent}\bfseries, commentstyle=\color{muted}\itshape,
  frame=single, rulecolor=\color{rule}, framesep=4pt, xleftmargin=4pt,
  xrightmargin=4pt, columns=fullflexible, keepspaces=true, breaklines=true,
  showstringspaces=false, upquote=true, aboveskip=0.6em, belowskip=0.6em}
\newcommand{\cmd}[1]{\texttt{#1}}
\newcommand{\file}[1]{\texttt{#1}}

\newcommand{\doctitle}{Design Specification}

\begin{document}
\makedoctitle{Architecture, arithmetic contract, and verification of the int8 MNIST accelerator}{2026-10-16}

\begin{note}[How to use this document]
This is the \emph{reference}: what the hardware does and why. If the RTL and this
document disagree, one of them is a bug, so fix whichever is wrong in the same
pull request. Day-to-day status lives in \file{progress.tex}; installing the
tools lives in \file{setup.tex}.
\end{note}

\tableofcontents

% ============================================================================
\section{Goal and scope}

We are building a hardware inference engine for a small neural network that
classifies handwritten digits (MNIST), running on a Digilent Basys~3 board
(AMD Artix-7 \texttt{xc7a35tcpg236-1}, 100\,MHz clock).

\begin{itemize}
  \item \textbf{Network:} a multilayer perceptron, $784 \rightarrow 32 \rightarrow 10$,
        ReLU on the hidden layer. Trained in PyTorch, then quantized to 8-bit integers.
  \item \textbf{Hardware does inference only.} All training happens offline in Python.
        The FPGA runs the forward pass using fixed, pre-computed weights.
  \item \textbf{Correctness standard:} the hardware must match the Python integer
        model \emph{bit for bit}, not approximately. Every intermediate value is
        checked against a golden reference (Section~\ref{sec:verif}).
  \item \textbf{Deliverable:} classify MNIST test images sent from a PC over UART
        and report accuracy measured on real hardware.
\end{itemize}

\textbf{Explicitly out of scope} for the first version: training on the FPGA,
camera or other live input, convolutional layers, and external DDR memory (all
weights fit on-chip, so none is needed).

% ============================================================================
\section{Model and quantization}

\subsection{Training}
\file{train\_mlp.py} trains the float32 reference network.

\begin{center}
\begin{tabular}{@{}ll@{}}
\toprule
Architecture  & $784 \rightarrow 32$ (ReLU) $\rightarrow 10$, with biases \\
Optimizer     & Adam, learning rate $10^{-3}$, batch 128, 8 epochs, seed 0 \\
Input scaling & pixels divided by 255 to $[0,1]$; \emph{no} mean/std normalization \\
\bottomrule
\end{tabular}
\end{center}

Skipping mean/std normalization is deliberate: it keeps every input
non-negative, so the hardware needs no signed input path. No softmax is
implemented, because $\arg\max$ gives the same answer with or without it.

\subsection{Integer scheme}
\file{quantize\_export.py} converts the model to integers and \textbf{is the
hardware specification}: whatever arithmetic it performs, the RTL must reproduce
exactly. Quantization is symmetric and per-tensor (one scale per tensor rather
than per neuron, which costs a little accuracy and saves a lot of hardware).

\begin{align*}
x_q &= \operatorname{round}(127\,x) \in [0, 127]
  && \text{input pixel (int8)}\\
W_q &= \operatorname{round}(W / s_W),\quad s_W = \max|W| / 127
  && \text{weights (int8)}\\
b_q &= \operatorname{round}\!\big(b / (s_x s_W)\big)
  && \text{biases (int32, accumulator scale)}\\[0.4em]
\mathit{acc}_1 &= \max\!\big(0,\; W_{1q}\,x_q + b_{1q}\big)
  && \text{layer 1 + ReLU (int32)}\\
h_q &= \min\!\big(127,\; (\mathit{acc}_1 + 2^{S-1}) \gg S\big),\quad S = 12
  && \text{requantize: round half up, clamp}\\
\mathit{acc}_2 &= W_{2q}\,h_q + b_{2q}
  && \text{layer 2 (int32)}\\
\hat{y} &= \arg\max(\mathit{acc}_2)
  && \text{predicted digit}
\end{align*}

The shift amount $S$ (\texttt{SHIFT1}) is the smallest shift that keeps the
largest post-ReLU layer-1 value inside int8 range, measured on 2{,}000
\emph{training} images. Calibrating on training images rather than test images
keeps test data out of the design, so the reported accuracy is honest.

\subsection{Measured results}
\begin{center}
\begin{tabular}{@{}lr@{}}
\toprule
float32 test accuracy (10{,}000 images)        & 95.86\% \\
int8 test accuracy                              & 95.86\% \\
Accuracy change from quantization               & $+0.00$ percentage points \\
Same prediction as float32                     & 99.74\% of images \\
\midrule
\texttt{SHIFT1}                                 & 12 (rounding on) \\
Largest post-ReLU layer-1 value (calibration)   & 294{,}415 \\
Accumulator headroom vs.\ int32                 & $7{,}294\times$ \\
\midrule
Weight storage, layer 1 / layer 2               & 25{,}088 B / 320 B (24.8\,KiB total) \\
\bottomrule
\end{tabular}
\end{center}

\begin{warn}[Weights are frozen]
Training is not bit-reproducible across machines. The committed files in
\file{rtl/} are the canonical weights; do not re-run \file{train\_mlp.py} unless
the team agrees to re-verify everything. \file{quantize\_export.py} is
deterministic given \file{mlp\_fp32.npz}, but that file is not in git (see the
open issue in \file{progress.tex}).
\end{warn}

\subsection{Generated files (\file{rtl/})}
\begin{center}
\begin{tabularx}{\linewidth}{@{}lX@{}}
\toprule
File & Contents \\
\midrule
\file{params.vh}        & Verilog \texttt{`define}s: sizes, \texttt{ACC\_W}, \texttt{SHIFT1}, \texttt{ROUND\_SHIFT}. Single source of truth; never edit by hand. \\
\file{w1.mem}           & Layer-1 weights, \textbf{neuron-major}: 784 weights for neuron 0, then neuron 1, \ldots{} One int8 per line. Used by \file{tb\_mac.v}. \\
\file{w1\_t.mem}        & Layer-1 weights, \textbf{pixel-major and packed}: 784 lines, each a 256-bit word holding pixel $i$'s weight for all 32 neurons. Neuron $n$ is bits $[8n{+}7{:}8n]$, so neuron 0 is the rightmost two hex digits. Read by \file{layer1.v}. \\
\file{b1.mem}, \file{b2.mem} & Biases, int32, two's complement hex. \\
\file{w2.mem}           & Layer-2 weights, neuron-major (10 $\times$ 32). \\
\file{test\_images.mem}, \file{test\_labels.txt} & First 100 MNIST test images (int8) and labels, for the on-board run. \\
\file{golden.txt}       & For the first 64 test images: label, prediction, and every intermediate value (\texttt{acc1}, \texttt{hq}, \texttt{acc2}). The testbenches compare against this. \\
\bottomrule
\end{tabularx}
\end{center}

% ============================================================================
\section{Hardware architecture}

\subsection{Block diagram}
\begin{center}
\begin{tikzpicture}[
  blk/.style={draw=accent, fill=accent!6, rounded corners=2pt, minimum height=1.1cm,
              align=center, font=\small, inner sep=5pt},
  plan/.style={blk, draw=todoc, fill=todoc!6, dashed},
  mem/.style={draw=muted, fill=codebg, minimum height=0.8cm, align=center, font=\footnotesize, inner sep=4pt},
  arr/.style={-{Stealth[length=2.2mm]}, thick, color=muted},
  lbl/.style={font=\scriptsize, color=muted}]
  \node[plan] (src) {Image buffer\\\scriptsize 784 B BRAM\\\scriptsize (from UART)};
  \node[blk, right=1.5cm of src, minimum width=3.6cm] (l1) {\textbf{Layer 1}\\\scriptsize 32 MACs in parallel\\\scriptsize ReLU + requantize};
  \node[plan, right=1.1cm of l1] (l2) {Layer 2\\\scriptsize 10 MACs};
  \node[plan, right=1.5cm of l2] (am) {argmax};
  \node[mem, above=0.55cm of l1] (wr) {weight ROM \file{w1\_t.mem}\\784 $\times$ 256 bit};
  \draw[arr] (src) -- node[above, lbl]{pixel} node[below, lbl]{1/clk} (l1);
  \draw[arr] (wr) -- (l1);
  \draw[arr] (l1) -- node[above, lbl]{$h[32]$} node[below, lbl]{int8} (l2);
  \draw[arr] (l2) -- node[above, lbl]{10 scores} node[below, lbl]{int32} (am);
  \draw[arr] (am.east) -- ++(0.7,0) node[right, lbl]{digit};
  \node[font=\scriptsize, color=muted, below=0.45cm of l1.south east, anchor=north] {solid = built and verified \quad dashed = planned};
\end{tikzpicture}
\end{center}

\subsection{MAC unit (\file{src/mac.v})}
One multiply-accumulate per clock: $\mathit{acc} \leftarrow \mathit{acc} + a \cdot b$.

\begin{center}
\begin{tabularx}{\linewidth}{@{}llX@{}}
\toprule
Signal & Width & Behavior \\
\midrule
\texttt{load} & 1 & At the clock edge, $\mathit{acc} \leftarrow \texttt{init}$ (the neuron's bias). \textbf{Has priority over \texttt{en}.} \\
\texttt{en}   & 1 & At the clock edge, $\mathit{acc} \leftarrow \mathit{acc} + a\cdot b$. Otherwise holds. \\
\texttt{a}, \texttt{b} & 8 signed & Activation and weight. The 16-bit product cannot overflow. \\
\texttt{init}, \texttt{acc} & 32 signed & Bias in, running sum out. \\
\bottomrule
\end{tabularx}
\end{center}

The \texttt{(* use\_dsp = "yes" *)} attribute maps each MAC onto one DSP48E1 slice.
The Basys~3 has 90, which is the real limit on how many MACs can run in parallel.

\subsection{Layer 1 (\file{src/layer1.v})}
All 32 hidden neurons are computed at once. Each clock, \emph{one} pixel is
broadcast to all 32 MACs, and one 256-bit read of the weight ROM gives every MAC
its own weight for that pixel. Layer 1 therefore takes 784 accumulate cycles,
regardless of the number of neurons (up to the DSP limit).

\begin{center}
\begin{tabularx}{\linewidth}{@{}llX@{}}
\toprule
Port & Dir/width & Meaning \\
\midrule
\texttt{start}    & in, 1   & Pulse for one cycle while idle to process an image. Ignored while busy. \\
\texttt{busy}     & out, 1  & High from the cycle after \texttt{start} until \texttt{done}. \\
\texttt{done}     & out, 1  & One-cycle pulse; \texttt{h\_flat} is valid from then until the next start. \\
\texttt{pix\_addr}& out, 10 & Index (0--783) of the pixel wanted next. \\
\texttt{pix\_data}& in, 8   & That pixel, \textbf{valid one cycle after \texttt{pix\_addr}} (synchronous BRAM timing). \\
\texttt{h\_flat}  & out, 256& 32 int8 activations; neuron $n$ in bits $[8n{+}7{:}8n]$. \\
\texttt{acc\_flat}& out, 1024& Raw int32 accumulators (pre-ReLU), exposed for verification. \\
\bottomrule
\end{tabularx}
\end{center}

\textbf{Timing.} Address $k$ is presented during the cycle after clock edge $k$.
\begin{center}
\begin{tabular}{@{}ll@{}}
\toprule
Edge & Event \\
\midrule
0       & \texttt{start} sampled; FSM leaves idle \\
1       & all MACs load their bias \\
$k+1$   & pixel $k$ and weight word $k$ arrive from memory \\
$k+2$   & MACs accumulate pixel $k$ \quad ($k = 0 \ldots 783$, so edges 2--785) \\
786     & requantized outputs registered; \texttt{done} goes high \\
\bottomrule
\end{tabular}
\end{center}
\textbf{Latency: 786 cycles = 7.86\,\textmu s at 100\,MHz.}

\subsection{Resource estimate}
Synthesized for Artix-7 with the open-source Yosys flow; Vivado's numbers will
be recorded here once the full design is implemented.

\begin{center}
\begin{tabular}{@{}lrrr@{}}
\toprule
Resource & Layer 1 uses & Basys~3 has & Share \\
\midrule
DSP48E1 slices        & 32     & 90      & 36\% \\
Block RAM (RAMB18)    & 15     & 100     & 15\% \\
LUTs                  & 1{,}657 & 20{,}800 & 8\% \\
Flip-flops            & 1{,}299 & 41{,}600 & 3\% \\
\bottomrule
\end{tabular}
\end{center}

% ============================================================================
\section{Verification}\label{sec:verif}

\subsection{Strategy}
\begin{enumerate}
  \item Python produces \textbf{golden vectors}: the exact integer value of every
        intermediate signal for 64 real test images (\file{golden.txt}).
  \item Each testbench runs the RTL on the same inputs and compares every value.
        Python and Verilog are independent implementations, so agreement means
        the hardware matches the model.
  \item Every testbench is \textbf{self-checking}: it prints
        \texttt{ALL TESTS PASSED} or names the exact failing value. Nobody has to
        read waveforms to know whether it worked.
  \item Every testbench is checked by \textbf{mutation testing}: we plant
        realistic bugs in the RTL and confirm the testbench catches each one.
        A testbench that passes is only meaningful if it can also fail.
\end{enumerate}

\subsection{Testbench results}
\begin{center}
\begin{tabularx}{\linewidth}{@{}lXrl@{}}
\toprule
Testbench & What it checks & Checks & Status \\
\midrule
\file{tb\_mac.v} & 10 directed edge cases (sign, extremes, hold, load priority);
  2{,}000 random dot products; 8 images $\times$ 32 neurons vs.\ golden & 2{,}266 & \done \\
\file{tb\_layer1.v} & 64 golden images $\times$ 32 neurons, both \texttt{acc1} and
  \texttt{hq}; 3 synthetic images; latency, \texttt{busy}, \texttt{done}, and a
  stray \texttt{start} on every run & 4{,}259 & \done \\
\bottomrule
\end{tabularx}
\end{center}
Both pass in Vivado's simulator (xsim) and in Icarus Verilog
(\file{scripts/run\_sims.sh}).

\subsection{Why \file{tb\_layer1} needs synthetic images}
Mutation testing showed that real MNIST digits have blind spots. Their border
pixels are always black, so dropping the first or last pixel changes no result.
No real image pushes a hidden neuron past 127, so a missing clamp is
invisible. Three made-up images close those gaps: a \emph{ramp} (every pixel
nonzero and different), \emph{all-127}, and a \emph{hot} image that lights every
pixel where the strongest neuron has a positive weight. Their expected values
are computed from \file{w1.mem} (neuron-major), not the \file{w1\_t.mem} the
hardware reads, so a packing error in \file{w1\_t.mem} is caught too. The test
prints how many times ReLU (41) and the clamp (2) were exercised, as proof those
paths were tested.

\subsection{Mutation testing results}
\textbf{\file{tb\_layer1.v}}: number of failing checks per planted bug. Every bug is caught.
\begin{center}
\begin{tabular}{@{}lrrrrl@{}}
\toprule
Planted bug & acc1 & hq & synthetic & control & \\
\midrule
MAC enable driven a cycle early (pipeline slip)   & 0    & 0    & 52 & 0  & caught \\
Weight memory one address behind pixels           & 2048 & 2048 & 96 & 0  & caught \\
Truncate instead of round when requantizing       & 0    & 705  & 24 & 0  & caught \\
Weight bytes in reversed order                    & 1875 & 1811 & 96 & 0  & caught \\
Bias not reloaded between images                  & 2048 & 2048 & 96 & 0  & caught \\
ReLU missing before requantize                    & 0    & 583  & 41 & 0  & caught \\
Last pixel skipped                                & 0    & 0    & 52 & 67 & caught \\
First pixel skipped                               & 0    & 0    & 87 & 67 & caught \\
Clamp to 127 missing                              & 0    & 0    & 2  & 0  & caught \\
Shift off by one                                  & 0    & 1437 & 53 & 0  & caught \\
\bottomrule
\end{tabular}
\end{center}
Note the first row: a pipeline-timing bug that passes \emph{every} real image and
is caught only by the synthetic stage.

\textbf{\file{tb\_mac.v}}: unsigned multiply (2{,}261 failures), \texttt{en}
overriding \texttt{load} (caught only by the directed test), and a 16-bit
truncated accumulator (caught only by the random and golden tests). Each tier
catches bugs the others miss.

% ============================================================================
\section{Planned work and open design questions}

\begin{itemize}
  \item \textbf{Layer 2} ($32 \rightarrow 10$): same broadcast pattern with 10 MACs and
        32 accumulate cycles. Needs a pixel-major \file{w2\_t.mem} (32 words $\times$
        80 bits) from the exporter. No requantize: \texttt{acc2} goes straight to argmax.
  \item \textbf{argmax}: must break ties like \texttt{numpy.argmax}: the
        \emph{lowest} index wins. A strict ``greater than'' compare scanning from
        index 0 does this. Verify with a synthetic tie.
  \item \textbf{Image buffer + UART}: 784-byte BRAM written by a UART receiver;
        prediction sent back by a UART transmitter. Protocol not yet defined.
        The host can send raw MNIST bytes and the board converts each with
        $\texttt{pixel} \gg 1$. This was checked: for all 256 byte values it gives
        exactly the same int8 as \file{quantize\_input} in Python.
  \item \textbf{Timing closure} at 100\,MHz once the full design is assembled.
\end{itemize}

% ============================================================================
\section{Repository layout}
\begin{center}
\begin{tabularx}{\linewidth}{@{}lX@{}}
\toprule
\file{src/}     & Synthesizable Verilog (\file{mac.v}, \file{layer1.v}, \ldots) \\
\file{sim/}     & Self-checking testbenches (\file{tb\_*.v}) \\
\file{rtl/}     & Generated weights, golden vectors, \file{params.vh}. Committed; canonical. \\
\file{scripts/} & \file{create\_project.tcl} (rebuilds the Vivado project),
                  \file{run\_sims.sh} (runs every testbench in Icarus) \\
\file{docs/}    & This document, \file{progress.tex}, \file{setup.tex}, and their PDFs \\
\file{vivado/}  & Generated Vivado project. \textbf{Not committed}; rebuilt per machine. \\
\file{*.py}     & \file{train\_mlp.py} (training), \file{quantize\_export.py} (the spec) \\
\bottomrule
\end{tabularx}
\end{center}

\end{document}
